% 課題別の逆引きガイドの流れ図。
% reader_guide_flow.pdf（PDF用）と reader_guide_flow.png（Web・EPUB用）の元データです。
% 生成: scripts/build-reader-guide-flow.sh
\documentclass[uplatex,dvipdfmx,class=ujarticle,border=1mm]{standalone}
\usepackage[deluxe,uplatex]{otf}
\usepackage[T1]{fontenc}
\usepackage[noalphabet,unicode]{pxchfon}
\setgothicfont{BIZUDGothic-Regular.ttf}
\setboldgothicfont{BIZUDGothic-Bold.ttf}
\renewcommand{\sfdefault}{phv}
\renewcommand\familydefault{\sfdefault}
\renewcommand\kanjifamilydefault{\gtdefault}
\usepackage{tikz}
\usetikzlibrary{positioning,arrows.meta}
\definecolor{ink}{HTML}{2B2B2B}
\definecolor{panel}{HTML}{F2F4F7}
\definecolor{sub}{HTML}{55606E}

\newcommand{\colw}{33mm}
\newcommand{\colgap}{36mm}

\begin{document}
\kanjiskip=.05zw plus.5pt minus.5pt\xkanjiskip=0pt
\begin{tikzpicture}[
  font=\sffamily\gtfamily,
  line/.style={draw=ink, line width=0.6pt},
  arrow/.style={line, -{Stealth[length=2mm]}},
  measure/.style={draw=ink, line width=0.8pt, rounded corners=1.5mm, fill=white,
    align=center, inner sep=1.8mm, text width=120mm},
  symptom/.style={align=center, font=\small\bfseries\gtfamily, text=ink},
  section/.style={fill=ink, text=white, font=\bfseries\gtfamily, minimum height=7mm,
    text width=\colw, align=center, inner sep=0pt},
  leaf/.style={fill=white, draw=ink!25, line width=0.4pt, text width=\colw-4mm,
    align=left, inner xsep=2mm, inner ysep=1.6mm, font=\footnotesize\gtfamily},
]
% 計測
\node[measure] (measure) at (1.5*\colgap, 0) {%
  {\large\bfseries 計測する}\\[0.8mm]
  {\small 進め方・基本ツール：前作 第1章・第3章\qquad
  本書の計測手段：「まずは計測から」}};

% 症状
\node[symptom] (s1) at (0, -17mm) {アプリ容量が大きい};
\node[symptom] (s2) at (\colgap, -17mm) {メモリ超過・クラッシュ};
\node[symptom] (s3) at (2.5*\colgap, -17mm) {処理落ち・スパイク};

\draw[line] (measure.south) -- ++(0, -4mm) coordinate (bus);
\draw[line] (0, 0 |- bus) -- (2.5*\colgap, 0 |- bus);
\foreach \s in {s1, s2, s3} \draw[arrow] (\s |- bus) -- (\s.north);

% 課題（本書の節）
\foreach \x/\name/\label in {0/c1/容量, 1/c2/メモリ, 2/c3/CPU, 3/c4/GPU}
  \node[section] (\name) at (\x*\colgap, -28mm) {\label};
\draw[arrow] (s1.south) -- (c1.north);
\draw[arrow] (s2.south) -- (c2.north);
\draw[line] (s3.south) -- ++(0, -1.6mm) coordinate (split);
\draw[line] (c3 |- split) -- (c4 |- split);
\draw[arrow] (c3 |- split) -- (c3.north);
\draw[arrow] (c4 |- split) -- (c4.north);

% 症状ごとの章
% \leafitem{上のノード}{名前}{症状}{章}
% \ch{部-章}{短縮した章タイトル}：症状と章、章と章の間を1行空け、2行目以降は番号の後ろにそろえる
\newcommand{\ch}[2]{\par
  \vspace{\baselineskip}
  \hangindent=2.2em\hangafter=1
  \noindent\makebox[2.2em][l]{#1}#2}
\newcommand{\leafitem}[4]{%
  \node[leaf, below=1.6mm of #1] (#2) {{\bfseries #3}\par{\color{sub}#4\par}};}

\leafitem{c1}{c1a}{アセットが増え続ける}{\ch{3-1}{AssetBundleのライフタイム管理}\ch{3-2}{アセット肥大化の調査}}
\leafitem{c1a}{c1b}{コードが大きい}{\ch{3-4}{Code Strippingによるコードの削減}\ch{3-5}{LTOによるコードの最適化}}

\leafitem{c2}{c2a}{上限で落ちる}{\ch{3-6}{iOSのメモリ上限}}
\leafitem{c2a}{c2b}{解放されない}{\ch{2-5}{UniTaskのライフタイム管理}\ch{3-1}{AssetBundleのライフタイム管理}}
\leafitem{c2b}{c2c}{描画のメモリ}{\ch{1-1}{Render Graphのタイルメモリ最適化}\ch{1-5}{RSUVによる描画の最適化}\ch{2-3}{VATによるボーンアニメーションの軽量化}}

\leafitem{c3}{c3a}{キャラクターが多い}{\ch{2-1}{キャラクターアニメーションのチューニング}\ch{2-2}{Animatorプーリングの初期化}\ch{2-3}{VATによるボーンアニメーションの軽量化}}
\leafitem{c3a}{c3b}{単発のスパイク}{\ch{2-4}{ParticleSystemのSimulateのスパイク}}
\leafitem{c3b}{c3c}{描画の登録・UI}{\ch{1-2}{UnsafePassによる描画登録の軽量化}\ch{2-6}{UIシステムの使い分け}}
\leafitem{c3c}{c3d}{C\#コード全体}{\ch{3-5}{LTOによるコードの最適化}}

\leafitem{c4}{c4a}{メモリ帯域}{\ch{1-1}{Render Graphのタイルメモリ最適化}\ch{1-3}{FinalBlitPassの省略}}
\leafitem{c4a}{c4b}{シェーダーの演算}{\ch{1-4}{シェーダーの超越関数の軽量化}}
\leafitem{c4b}{c4c}{初回描画のスタッター}{\ch{1-6}{シェーダーウォームアップ}}
\leafitem{c4c}{c4d}{大量描画・端末差}{\ch{1-5}{RSUVによる描画の最適化}\ch{3-3}{端末別の描画設定}}
\end{tikzpicture}
\end{document}
